\documentclass[11pt,a4paper]{article}
\usepackage{fontspec}
\usepackage{polyglossia}
\setmainlanguage{russian}
\setotherlanguage{english}
\setmainfont{Times New Roman}
\setmonofont{DejaVu Sans Mono}
\usepackage{amsmath,amssymb}
\usepackage{geometry}
\usepackage{booktabs}
\usepackage{array}
\usepackage{microtype}
\usepackage[hidelinks]{hyperref}
\geometry{margin=22mm}
\setlength{\parindent}{1.25em}
\setlength{\parskip}{0.35em}
\setlength{\emergencystretch}{3em}
\newcommand{\code}[1]{\path{#1}}

\title{Отчёт по трём пилотным экспериментам}
\author{}
\date{25 сентября 2026 г.}

\begin{document}
\maketitle

\section*{Короткий вывод}

Наиболее содержательный результат дала плотная low-rank matrix factorization с momentum без внешнего forcing. При $\beta=1$ полные матрицы факторов сами порождают длительную эндогенную динамику, а оценка по скалярным логам возрастает при переходе от задач ранга 1 к рангам 2 и 4. Контроль с $\beta=0.9$ даёт затухающую траекторию и менее устойчивую оценку.

Эксперимент с двухслойной линейной сетью пока не подтвердил переносимость метода на нейросетевую оптимизацию: оценка почти не зависела от ранга учителя и не прошла проверку идентифицируемости. Автоматический выбор лага после смены постановки требует отдельного пересчёта.

\section{Плотная matrix factorization с momentum}

\subsection{Постановка}

Мы проверяем метод на эндогенной оптимизационной траектории, для которой известна контрольная сложность задачи. Учится обычная плотная факторизация:

\[
    \min_{U,V}\frac12\left\|UV^{\mathsf T}-M_\star\right\|_F^2,
    \qquad \operatorname{rank}(M_\star)=r.
\]

Целевая матрица строится как

\[
    M_\star=P\operatorname{diag}(\sigma_1,\ldots,\sigma_r)Q^{\mathsf T},
\]

где $P$ и $Q$ --- случайные ортогональные матрицы, а $\sigma_i$ --- ненулевые сингулярные значения. Поэтому $M_\star$ плотная и не диагональна в координатах обучения. Обучаются полные плотные матрицы $U$ и $V$ ширины 8; диагонального ограничения на факторы нет. Варьируется ранг цели $r\in\{1,2,4\}$.

Здесь $r$ --- известный контроль сложности задачи. Мы не утверждаем, что estimator обязан вернуть ровно $r$: у плотной факторизации есть дополнительные параметрические направления и симметрии. Проверяется, реагирует ли оценка на изменение ранга и сохраняет ли устойчивую динамику.

Обучение full-batch, без шума, случайных мини-батчей и внешнего forcing. Используется обновление

\[
    v_{t+1}=\beta v_t+\nabla L(\theta_t),
    \qquad
    \theta_{t+1}=\theta_t-\eta v_{t+1}.
\]

Для основного эксперимента $\eta=0.02$, $\beta=1$, 4000 шагов и два random seed. Для контроля используется $\beta=0.9$, при котором движение затухает. Сохраняются loss, нормы параметров и velocity, а также случайная линейная проекция полной траектории.

Оценка выполняется с embedding dimension до 6, $k=4$, Theiler window в embedding-режиме и локальным окном длиной 500 наблюдений. Проверяются лаги $\tau=4$, $\tau=8$ и лаг по автокорреляции.

\subsection{Что измеряется}

Контрольный ориентир --- ранг целевой матрицы $r$. Для каждого scalar log вычисляются $MG$, идентифицируемость, число пересечений тренда и флаг вырожденного окна.

Оценка считается пригодной только вместе с диагностикой, а не по одному числу $MG$.

\subsection{Результат}

В recurrent-режиме ни один из проверенных вариантов не помечен как вырожденный. Медианные значения $MG$:

\begin{center}
\begin{tabular}{>{\raggedleft\arraybackslash}p{0.29\linewidth} >{\centering\arraybackslash}p{0.27\linewidth} >{\centering\arraybackslash}p{0.27\linewidth}}
\toprule
Ранг цели $r$ & $MG$ по loss & $MG$ по velocity \\
\midrule
1 & 3.5 & 3.5 \\
2 & 5.5 & 5.5 \\
4 & 5.9 & 5.9 \\
\bottomrule
\end{tabular}
\end{center}

Оценка возрастает при переходе от ранга 1 к рангу 2 и далее к рангу 4. Абсолютное значение не совпадает с $r$, что ожидаемо: размерность траектории факторов определяется не только рангом цели, но и параметризацией, momentum и выбранным scalar observable.

В контроле с $\beta=0.9$ движение затухает, а оценки становятся заметно выше: медианный $MG$ составляет примерно 17--18 против 4--6 в recurrent-режиме. Текущая диагностика не всегда помечает такие окна как вырожденные, что показывает слабость reject-механизма для затухающих плотных траекторий.

\subsection{Вывод}

Плотная factorization является более содержательной controlled validation, чем диагональная версия: координаты мод не заданы заранее, а факторы обучаются как обычные плотные матрицы. Результат показывает реакцию scalar-log protocol на увеличение ранга задачи на эндогенной траектории без внешнего forcing.

Это не точное восстановление active dimension по рангу матрицы. Результат остаётся пилотным: два seed, искусственная задача и недостаточно строгий reject-механизм для затухающих траекторий.

\section{Teacher--student с двухслойной линейной сетью}

\subsection{Постановка}

Учится двухслойная линейная сеть $f(x)=W_2W_1x$ размерности 8--8--8. Данные генерируются фиксированным линейным учителем с $r\in\{1,2,4\}$ ненулевыми сингулярными значениями. Обучение full-batch с momentum $\beta=1$, $\eta=0.02$, 3000 шагов и двумя seed; обучаются все элементы $W_1$ и $W_2$.

\subsection{Результат}

Диагностика не объявляет окна вырожденными, но проверка идентифицируемости неудовлетворительна: отношение оценок обычно находится в диапазоне 1.7--2.1. Медианный $MG$ почти не зависит от ранга учителя:

\begin{center}
\begin{tabular}{>{\raggedleft\arraybackslash}p{0.29\linewidth} >{\centering\arraybackslash}p{0.27\linewidth} >{\centering\arraybackslash}p{0.27\linewidth}}
\toprule
Ранг учителя & $MG$ по loss & $MG$ по velocity \\
\midrule
1 & 6.1 & 6.0 \\
2 & 6.0 & 6.1 \\
4 & 6.2 & 5.9 \\
\bottomrule
\end{tabular}
\end{center}

Этот запуск не подтверждает метод на нейросетевой оптимизации. Вероятные причины --- недостаточно выраженная recurrent-динамика, параметры learning rate/momentum и симметрии двухслойной факторизации.

\section{Автоматический выбор лага и окна}

Проверялись фиксированные лаги $\tau=4$, $\tau=8$, лаг по автокорреляции и пилотный выбор лага по доминирующему периоду periodogram scalar log.

После изменения первой постановки автоматический выбор лага ещё не пересчитан для плотной factorization. Поэтому старые числа periodogram-пилота не используются как результат новой controlled validation.

Для новой постановки нужно выбирать устойчивый диапазон лагов, проверять стабильность $MG$ по этому диапазону и иметь явный исход «режим не удалось идентифицировать».

\section*{Итоговая оценка}

\begin{enumerate}
  \item Плотная matrix factorization + momentum --- наиболее содержательная новая controlled validation. Она показывает рост оценки при увеличении ранга на эндогенной recurrent-траектории, но выявляет слабость текущего reject-механизма для затухающих траекторий.
  \item Teacher--student --- отрицательный результат для выбранных параметров и мотивация для более тщательно настроенного neural-network эксперимента.
  \item Automatic lag selection --- открытая часть эксперимента; старые результаты диагональной постановки нельзя переносить на плотную factorization без пересчёта.
\end{enumerate}

Данные экспериментов находятся в \code{raw\_scores.csv} и \code{summary.csv}.

\end{document}
